% Generated draft: requires mathematical and source review.
\documentclass[11pt]{article}
\usepackage[margin=25mm]{geometry}
\usepackage{fontspec}
\setmainfont{Latin Modern Roman}
\usepackage{amsmath,amssymb,mathtools,unicode-math}
\setmathfont{Latin Modern Math}
\usepackage{enumitem,needspace,xurl,hyperref}
\setlength{\parindent}{0pt}
\setlength{\parskip}{5pt}
\setlength{\emergencystretch}{3em}
\newcommand{\N}{\mathbb{N}}
\newcommand{\Z}{\mathbb{Z}}
\newcommand{\Q}{\mathbb{Q}}
\newcommand{\R}{\mathbb{R}}
\newcommand{\C}{\mathbb{C}}
\newcommand{\F}{\mathbb{F}}
\newcommand{\A}{\mathbb{A}}
\newcommand{\PP}{\mathbb{P}}
\newcommand{\E}{\mathbb{E}}
\newcommand{\eps}{\varepsilon}
\newcommand{\dd}{\,\mathrm{d}}
\begin{document}
\section*{UnsolvedMath ID: 11300001}
\textbf{Source problem code:} catalogue entry
\subsection{Definitions and mathematical statement}
For a topological space $X$, homogeneity means that for every $x,y\in X$ there is a homeomorphism $h:X\to X$ with $h(x)=y$. Under the usual metrizable-space convention, $X$ is an absolute neighborhood retract (ANR) if whenever it is embedded as a closed subspace of a metrizable space $Y$, some neighborhood $U$ of $X$ in $Y$ admits a retraction $U\to X$. Let $\dim X=n<\infty$ denote covering dimension.

The Bing--Borsuk question in the supplied source is whether every $n$-dimensional homogeneous ANR $X$ is a topological $n$-manifold. With the usual boundaryless manifold convention, the conclusion means that $X$ is Hausdorff and second countable and every point has an open neighborhood homeomorphic to an open subset of $\R^n$. The source does not explicitly specify the ANR category, separability, local compactness, connectedness, or compactness; the precise standard formulation should be checked before treating these conventions as definitive. No compactness restriction is silently imposed, and no conclusion is substituted merely about homology manifolds.

The related modified conjecture, mentioned only in the research summary, asks whether every finite-dimensional homogeneous ANR is a homology manifold. Its local integral homology condition in dimension $n$ is
\[
 H_k(X,X\setminus\{x\};\Z)\cong
 \begin{cases}\Z,&k=n,\\0,&k\ne n,\end{cases}
 \qquad x\in X,
\]
with the ambient regularity and homology-manifold conventions still needing specification. This is a weaker related target, not the topological-manifold conclusion of the main question.
\subsection{Short English statement}
Must a finite-dimensional homogeneous absolute neighborhood retract be a topological manifold? The supplied assessment reports the conjecture open, known in dimensions at most two, and accompanied by substantial structural results. It separately reports a March 2018 counterexample announcement by Bryant and Ferry as unpublished; this is retained as an attributed announcement rather than an established disproof.
\subsection{Sources}
\begin{itemize}
\item Local UnsolvedMath registry, permanent ID $11300001$, source code \texttt{catalogue entry}, attributed to R. H. Bing and Karol Borsuk; created 2026-07-29, updated 2026-08-06. Supplied status \texttt{partially\_solved}; classification \texttt{PARTIAL-PROGRESS}; difficulty L4, also research-suggested L4, while extraction background records default L3.
\item \emph{Topology problems from Wikipedia}, ``Bing--Borsuk conjecture,'' MediaWiki source revision 1366636767, \url{https://en.wikipedia.org/w/index.php?title=List_of_unsolved_problems_in_mathematics&oldid=1366636767#Topology}, accessed 2026-07-29. The background reports exhaustive reconciliation of a corrected current topology section against v1.1 and authoritative status sources and calls this a distinct still-open entry; those sources are not individually identified.
\item The research summary attributes the $n\leq2$ cases to classical Bing--Borsuk work and a dimension-three implication of the Poincar\'e conjecture to Jacobsche, describing the three-dimensional case as open. This is a reported implication, not a statement that the Poincar\'e conjecture itself is open. It reports that J. Bryant and S. Ferry announced a counterexample at the Spring Topology Conference in March 2018 but that it had not been published in the supplied assessment.
\item The summary describes the modified conjecture as open and names Bredon and Bryant for partial results involving finitely generated local homology groups. It mentions the Homogeneity and Resolution conjectures as companions and structural theory of homogeneous ANR compacta. Exact theorem hypotheses and bibliographic references are not supplied.
\item Source-status evidence separately remains \texttt{NEEDS\_REVIEW}, with a current primary-literature audit outstanding. Rights note \texttt{NEEDS\_REVIEW}: publicly accessible author/publisher copy, redistribution terms requiring review. Dataset attribution: Ulam AI / UnsolvedMath Contributors, CC BY 4.0.
\end{itemize}
\end{document}
